% ===================================================================== PREAMBULE COMMUN — Carnet d'entraînement Fiche 9
\documentclass[11pt,a4paper]{article}
\usepackage[utf8]{inputenc}
\usepackage[T1]{fontenc}
\usepackage{lmodern}
\usepackage[french]{babel}
\usepackage{amsmath,amssymb,mathtools}
\usepackage{icomma}
\usepackage{siunitx}
\sisetup{output-decimal-marker={,},per-mode=symbol}
\DeclareSIUnit{\molar}{M}
\DeclareSIUnit{\Molar}{mol\,L^{-1}}
\usepackage[version=4]{mhchem}
\usepackage{xcolor}
\usepackage{colortbl}
\usepackage{tikz}
\usepackage{pgfplots}
\usepackage[most]{tcolorbox}
\usepackage{enumitem}
\usepackage{array}
\usepackage{booktabs}
\usepackage{multicol}
\usepackage{fancyhdr}
\usepackage[hidelinks]{hyperref}
\usetikzlibrary{arrows.meta,positioning,calc,decorations.pathmorphing,
  decorations.markings,angles,quotes,patterns,backgrounds,
  shapes.geometric,shapes.misc,fadings,intersections,fit}
\pgfplotsset{compat=1.18}
\usepackage[a4paper,margin=2.2cm,headheight=26pt,headsep=10pt]{geometry}

% ---------- Palette ----------
\definecolor{primary}{RGB}{150,98,162}
\definecolor{primarydark}{RGB}{92,52,110}
\definecolor{defcol}{RGB}{0,114,178}
\definecolor{thmcol}{RGB}{0,140,120}
\definecolor{formulacol}{RGB}{198,93,9}
\definecolor{remarkcol}{RGB}{120,94,200}
\definecolor{attncol}{RGB}{200,40,55}
\definecolor{excol}{RGB}{0,135,90}
\definecolor{methcol}{RGB}{90,90,100}
\definecolor{cationcol}{RGB}{0,90,190}
\definecolor{anioncol}{RGB}{200,60,55}
\definecolor{cristalcol}{RGB}{110,105,120}
\definecolor{eaucol}{RGB}{120,180,220}
\definecolor{solcol}{RGB}{0,135,90}
\definecolor{kpscol}{RGB}{198,93,9}
\definecolor{qcol}{RGB}{120,94,200}
\definecolor{precipcol}{RGB}{110,70,40}
\definecolor{exercol}{RGB}{150,98,162}
\definecolor{corrcol}{RGB}{0,120,100}

% ---------- Macros ----------
\newcommand{\Kps}{K_{\mathrm{ps}}}
\newcommand{\Ce}[1]{\left[\,\ce{#1}\,\right]}

% ---------- Style des boîtes ----------
\tcbset{boxbase/.style={enhanced,breakable,boxrule=0.9pt,arc=2.5pt,
   left=3mm,right=3mm,top=2.3mm,bottom=2.3mm,
   fonttitle=\bfseries,coltitle=white,toptitle=0.6mm,bottomtitle=0.6mm}}

\newtcolorbox[auto counter]{definition}[1][]{%
  boxbase,colback=defcol!5,colframe=defcol,
  title={\textsc{Définition}~\thetcbcounter\ \ifx\relax#1\relax\relax\else\textmd{--- \itshape #1}\fi}}
\newtcolorbox[auto counter]{theoreme}[1][]{%
  boxbase,colback=thmcol!6,colframe=thmcol,
  title={\textsc{Loi / Propriété}~\thetcbcounter\ \ifx\relax#1\relax\relax\else\textmd{--- \itshape #1}\fi}}
\newtcolorbox[auto counter]{formule}[1][]{%
  boxbase,colback=formulacol!6,colframe=formulacol,
  title={\textsc{Formule}~\thetcbcounter\ \ifx\relax#1\relax\relax\else\textmd{--- \itshape #1}\fi}}
\newtcolorbox{remarque}[1][]{%
  boxbase,colback=remarkcol!6,colframe=remarkcol,
  title={\textsc{Remarque}\ifx\relax#1\relax\relax\else\textmd{ : \itshape #1}\fi}}
\newtcolorbox{attention}[1][]{%
  boxbase,colback=attncol!6,colframe=attncol,
  title={\textsc{Attention}\ifx\relax#1\relax\relax\else\textmd{ : \itshape #1}\fi}}
\newtcolorbox{exemple}[1][]{%
  boxbase,colback=excol!5,colframe=excol,
  title={\textsc{Exemple}\ifx\relax#1\relax\relax\else\textmd{ : \itshape #1}\fi}}
\newtcolorbox{methode}[1][]{%
  boxbase,colback=methcol!7,colframe=methcol,sharp corners,
  title={\textsc{Méthode}\ifx\relax#1\relax\relax\else\textmd{ : \itshape #1}\fi}}
\newtcolorbox{astuce}[1][]{%
  boxbase,colback=primary!7,colframe=primary,
  title={\textsc{Astuce}\ifx\relax#1\relax\relax\else\textmd{ : \itshape #1}\fi}}

\newtcolorbox[auto counter]{exercice}[1][]{%
  boxbase,colback=exercol!4,colframe=exercol,
  title={\textsc{Exercice}~\thetcbcounter\ \ifx\relax#1\relax\relax\else\textmd{--- \itshape #1}\fi}}
\newtcolorbox[auto counter]{correction}[1][]{%
  boxbase,colback=corrcol!4,colframe=corrcol,
  title={\textsc{Corrigé}~\thetcbcounter\ \ifx\relax#1\relax\relax\else\textmd{--- \itshape #1}\fi}}

% ---------- Environnement « où : » ----------
\newenvironment{ou}{%
  \par\smallskip\noindent\textbf{où :}\nopagebreak
  \begin{itemize}[leftmargin=1.8em,topsep=2pt,itemsep=1.5pt,parsep=0pt]}%
  {\end{itemize}}

% ---------- Styles TikZ ----------
\tikzset{
  cat/.style={circle,fill=cationcol,draw=cationcol!50!black,inner sep=0pt,
              minimum size=5mm,font=\bfseries\scriptsize,text=white},
  an/.style={circle,fill=anioncol,draw=anioncol!50!black,inner sep=0pt,
             minimum size=5mm,font=\bfseries\scriptsize,text=white},
  crx/.style={circle,fill=cristalcol,draw=black!50,inner sep=0pt,
              minimum size=5mm,font=\bfseries\scriptsize,text=white},
  ptn/.style={circle,fill=black,inner sep=1.1pt}}

% ---------- En-têtes ----------
\pagestyle{fancy}
\fancyhf{}
\fancyhead[L]{\small\textcolor{primarydark}{\textbf{Fiche 9} — Solubilité et produit de solubilité}}
\fancyhead[R]{\small\textcolor{primarydark}{\textsc{Carnet d'entraînement}}}
\fancyfoot[C]{\small\thepage}
\renewcommand{\headrulewidth}{0.8pt}
\renewcommand{\headrule}{\hbox to\headwidth{\color{primary}\leaders\hrule height \headrulewidth\hfill}}
\usepackage{titlesec}
\titleformat{\section}{\Large\bfseries\color{primarydark}}{\thesection}{0.6em}{}
\titleformat{\subsection}{\large\bfseries\color{primary}}{\thesubsection}{0.6em}{}

% ---------- Bandeau de titre ----------
% #1 = thème/étiquette   #2 = titre principal   #3 = sous-titre
\newcommand{\bandeau}[3]{%
\begin{center}
\begin{tikzpicture}
\node[rectangle,rounded corners=7pt,fill=primary,text=white,
      minimum width=15.6cm,minimum height=1.95cm,align=center,inner sep=8pt]
  {{\large\bfseries #1}\\[3pt]{\Large\bfseries #2}\\[3pt]{\small #3}};
\end{tikzpicture}
\end{center}
\vspace{2mm}}
\begin{document}
\bandeau{Thème 1 — Équation de dissolution \& $\Kps$}%
{Corrigé — Niveau Moyen}%
{Réponses détaillées et justifications}

\begin{correction}[phosphate de calcium]
\begin{enumerate}[label=\alph*),leftmargin=2em,topsep=2pt,itemsep=2pt]
  \item \textbf{Formule.} Il faut neutraliser : \ce{Ca^{2+}} ($+2$) et \ce{PO4^{3-}} ($-3$). Le plus
        petit commun multiple de 2 et 3 est 6 : il faut \textbf{3} \ce{Ca^{2+}} ($+6$) et \textbf{2}
        \ce{PO4^{3-}} ($-6$). D'où $\ce{Ca3(PO4)2}$.
  \item \textbf{Équation.} $\ce{Ca3(PO4)2(s) <=> 3 Ca^{2+}(aq) + 2 PO4^{3-}(aq)}$.
  \item \textbf{$\Kps$ et type.} $\Kps=\Ce{Ca^2+}^{3}\,\Ce{PO4^3-}^{2}$ ; type $\boldsymbol{3{:}2}$.
\end{enumerate}
\end{correction}

\begin{correction}[phosphate de baryum \ce{Ba3(PO4)2}]
\begin{enumerate}[label=\alph*),leftmargin=2em,topsep=2pt,itemsep=2pt]
  \item $\ce{Ba3(PO4)2(s) <=> 3 Ba^{2+}(aq) + 2 PO4^{3-}(aq)}$.
  \item $\Kps=\Ce{Ba^2+}^{3}\,\Ce{PO4^3-}^{2}$.
  \item Exposants : \textbf{3} sur \ce{Ba^{2+}} et \textbf{2} sur \ce{PO4^{3-}} (les coefficients de
        l'équation). Le groupement \ce{PO4^{3-}} est un \textbf{ion polyatomique stable} : il se
        déplace en bloc lors de la dissolution et ne se dissocie pas en \ce{P} et \ce{O}. On l'écrit
        donc comme une seule entité, avec sa charge $-3$.
\end{enumerate}
\end{correction}

\begin{correction}[lecture inverse]
\begin{enumerate}[label=\alph*),leftmargin=2em,topsep=2pt,itemsep=2pt]
  \item Les exposants indiquent les coefficients : 1 \ce{Ca^{2+}} et 2 \ce{F-}. Le sel neutre est donc
        $\ce{CaF2}$ (fluorure de calcium).
  \item $\ce{CaF2(s) <=> Ca^{2+}(aq) + 2 F-(aq)}$.
  \item Type $\boldsymbol{1{:}2}$.
\end{enumerate}
\emph{Méthode : l'exposant d'un ion dans le $\Kps$ = son coefficient dans l'équation = son indice dans
la formule.}
\end{correction}

\begin{correction}[tableau complété]
\begin{center}\renewcommand{\arraystretch}{1.5}
\begin{tabular}{|l|l|c|l|}
\hline
\rowcolor{primary!12}
\textbf{Sel} & \textbf{Équation} & \textbf{Type} & \textbf{$\Kps$}\\
\hline
\ce{BaSO4} & $\ce{BaSO4(s) <=> Ba^{2+} + SO4^{2-}}$ & $1{:}1$ & $\Ce{Ba^2+}\Ce{SO4^2-}$\\ \hline
\ce{PbI2}  & $\ce{PbI2(s) <=> Pb^{2+} + 2 I-}$      & $1{:}2$ & $\Ce{Pb^2+}\Ce{I-}^{2}$\\ \hline
\ce{Ag2S}  & $\ce{Ag2S(s) <=> 2 Ag+ + S^{2-}}$     & $2{:}1$ & $\Ce{Ag+}^{2}\Ce{S^2-}$\\ \hline
\ce{Cr(OH)3} & $\ce{Cr(OH)3(s) <=> Cr^{3+} + 3 OH-}$ & $1{:}3$ & $\Ce{Cr^3+}\Ce{OH-}^{3}$\\ \hline
\end{tabular}
\end{center}
Vérifie à chaque ligne : la somme des charges à droite est nulle, et l'exposant reproduit le coefficient.
\end{correction}

\begin{correction}[électrolyte ou non]
\begin{enumerate}[label=\alph*),leftmargin=2em,topsep=2pt,itemsep=2pt]
  \item \textbf{Électrolytes :} \ce{NaCl} et \ce{CaSO4} (composés ioniques $\to$ libèrent des ions). Le
        \textbf{glucose} \ce{C6H12O6} est \textbf{moléculaire} : il se dissout sans se dissocier en ions
        (non-électrolyte).
  \item $\ce{NaCl(s) <=> Na+(aq) + Cl-(aq)}$ \quad et \quad $\ce{CaSO4(s) <=> Ca^{2+}(aq) + SO4^{2-}(aq)}$.
  \item La notion de $\Kps$ n'a de sens que pour un \textbf{équilibre de dissolution ionique}. Le
        glucose reste sous forme de molécules \ce{C6H12O6(aq)} : il n'y a ni cation ni anion, donc
        \textbf{aucun produit de concentrations d'ions} à écrire — pas de $\Kps$.
\end{enumerate}
\end{correction}

\end{document}
